\begin{exercisechunk}
\item \textbf{A common target still requires coverage.}
\label[exercise]{ex:l7-covariate-coverage}
\exerciseprofile{2}{2}{2}{\exproof\quad\exconstruction\quad\exinterpretation}{%
Separate preservation of the pointwise target from identification of that
target by the training law.}
Let \(\cX=\{a,b\}\), \(\cY=\{0,1\}\), and let the loss be zero--one loss.
Suppose the training and deployment laws share the deterministic conditional
kernel
\[
 K(1\mid a)=0,
 \qquad
 K(1\mid b)=1.
\]
The model class contains every function from \(\cX\) to \(\cY\).
\begin{enumerate}[(a),beginpenalty=10000]
\item Prove directly from the two conditional risks that
\(f^\star(a)=0\), \(f^\star(b)=1\) minimizes population risk under every
input marginal on \(\cX\).
\item Let \(\mu_{\rm tr}(a)=1\) and \(\mu_{\rm dep}(b)=1\). Construct a
predictor with zero training risk and deployment risk one. Explain why the
training risk cannot distinguish this predictor from \(f^\star\).
\item Now assume
\(d\mu_{\rm dep}/d\mu_{\rm tr}\le \rho^{-1}\) for some
\(\rho\in(0,1]\). Prove the excess-risk bound in
\cref{prop:coverage} for this finite input space. Identify the step that fails
for the laws in part (b).
\end{enumerate}
\begin{hints}
For (c), write deployment excess risk as a sum over \(x\in\{a,b\}\) and
insert the density ratio into each summand.
\end{hints}
\end{exercisechunk}
